\subsection{Experimental Setup and Safeguards}
Every submission passes three gates: (i) an $\mathrm{Operator.equiv}$
circuit-identity assertion against the dense
$J^\dagger (U^{\otimes N}) J$ reference, (ii) a noiseless simulator dry-run
that must reproduce the exact analytic advantage, and (iii) a full dress rehearsal of the batch, including mitigation analysis, on the device
noise model reduced to the executed qubits. Job identifiers are persisted
before result polling, making runs crash-safe and recoverable. The execution
configuration common to all runs is summarised in Table~\ref{tab:hwconfig}.
% experiments/hardware_scaling.py; spec docs/superpowers/specs/2026-07-16-*.md

\begin{table}[t]
\centering
\caption{Hardware Execution Configuration}
\label{tab:hwconfig}
\footnotesize
\begin{tabular}{@{}ll@{}}
\toprule
\textbf{Parameter} & \textbf{Value}\\
\midrule
Device & ibm\_fez (IBM Heron r2)\\
Shots per circuit & 4096\\
$N{=}3$ validation chain & $[136,143,142]$\\
Scaling chain ($N{=}3,4,5$) & $[59,75,74,73,79]$\\
Extension chain ($N{=}3\ldots7$) & $[97,107,108,109,110,111,98]$\\
Readout mitigation & tensored confusion-matrix inversion\\
Zero-noise extrapolation & local cz folding, $\lambda\in\{1,3,5\}$\\
Batch sizes & 11-pub scaling, 17-pub extension,\\
 & 49-pub multi-topology\\
Software & qiskit 1.3.2 / qiskit-aer 0.14.2 /\\
 & qiskit-ibm-runtime 0.36.1\\
\bottomrule
\end{tabular}
\end{table}
